\section{AMC 10 10B 2004}
        \begin{enumerate}
        	\item (\textit{AMC 10 10B 2004}) Each row of the Misty Moon Amphitheater has $33$ seats. Rows $12$ through $22$ are reserved for a youth club. How many seats are reserved for this club?


 $\mathrm{(A) \ } 297 \qquad \mathrm{(B) \ } 330\qquad \mathrm{(C) \ } 363\qquad \mathrm{(D) \ } 396\qquad \mathrm{(E) \ } 726$


 
 \textbf{Answer: }C



\textbf{Solution}

There are $22-12+1=11$ rows of $33$ seats, giving $11\times 33=\boxed{\mathrm{(C)}\ 363}$ seats.

 


	\item (\textit{AMC 10 10B 2004}) How many two-digit positive integers have at least one $7$ as a digit?


 $\mathrm{(A) \ } 10 \qquad \mathrm{(B) \ } 18\qquad \mathrm{(C) \ } 19 \qquad \mathrm{(D) \ } 20\qquad \mathrm{(E) \ } 30$


 
 \textbf{Answer: }B



\textbf{Solution 1}

Ten numbers $(70,71,\dots,79)$ have $7$ as the tens digit. Nine numbers $(17,27,\dots,97)$ have it as the ones digit. Number $77$ is in both sets.

 Thus the result is $10+9-1=18 \Rightarrow$ $\boxed{\mathrm{(B)}\ 18}$.

 



\textbf{Solution 2}

We use complementary counting. The complement of having at least one $7$ as a digit is having no $7$s as a digit.

 We have $9$ digits to choose from for the first digit and $10$ digits for the second. This gives a total of $9 \times 10 = 90$ two-digit numbers.

 But since we cannot have $7$ as a digit, we have $8$ first digits and $9$ second digits to choose from.

 Thus there are $8 \times 9 = 72$ two-digit numbers without a $7$ as a digit.

 $90$ (The total number of two-digit numbers) $-  72$ (The number of two-digit numbers without a $7$) $= 18 \Rightarrow$ $\boxed{\mathrm{(B)}\ 18}$.

 


	\item (\textit{AMC 10 10B 2004}) At each basketball practice last week, Jenny made twice as many free throws as she made at the previous practice. At her fifth practice she made $48$ free throws. How many free throws did she make at the first practice?


 $\mathrm{(A) \ } 3 \qquad \mathrm{(B) \ } 6 \qquad \mathrm{(C) \ } 9 \qquad \mathrm{(D) \ } 12\qquad \mathrm{(E) \ } 15$


 
 \textbf{Answer: }A



\textbf{Solution}

At the fourth practice she made $48/2=24$ throws, at the third one it was $24/2=12$, then we get $12/2=6$ throws for the second practice, and finally $6/2=3\Rightarrow\boxed{\mathrm{(A)}\ 3}$ throws at the first one.

 


	\item (\textit{AMC 10 10B 2004}) A standard six-sided die is rolled, and $P$ is the product of the five numbers that are visible. What is the largest number that is certain to divide $P$?


 $\mathrm{(A) \ } 6 \qquad \mathrm{(B) \ } 12 \qquad \mathrm{(C) \ } 24 \qquad \mathrm{(D) \ } 144\qquad \mathrm{(E) \ } 720$


 
 \textbf{Answer: }B



\textbf{Solution 1}

The product of all six numbers is $6!=720$. The products of numbers that can be visible are $720/1$, $720/2$, ..., $720/6$.The answer to this problem is their greatest common divisor -- which is $720/L$, where $L$ is the least common multiple of $\{1,2,3,4,5,6\}$.Clearly $L=60$ and the answer is $720/60=\boxed{\mathrm{(B)}\ 12}$.

 



\textbf{Solution 2}

Clearly, $P$ cannot have a prime factor other than $2$, $3$ and $5$.

 We can not guarantee that the product will be divisible by $5$, as the number $5$ can end on the bottom.

 We can guarantee that the product will be divisible by $3$ (one of $3$ and $6$ will always be visible), but not by $3^2$.

 Finally, there are three even numbers, hence two of them are always visible and thus the product is divisible by $2^2$. This is the most we can guarantee, as when the $4$ is on the bottom side, the two visible even numbers are $2$ and $6$, and their product is not divisible by $2^3$.

 



\textbf{Solution 3}

The product P can be one of the following six numbers excluding the number that is hidden under, so we have:\begin{align*}2 \cdot 3 \cdot 4 \cdot 5 \cdot 6 = 2^4 \cdot 3^2 \cdot 5 \\1 \cdot 3 \cdot 4 \cdot 5 \cdot 6 = 2^3 \cdot 3^2 \cdot 5 \\1 \cdot 2 \cdot 4 \cdot 5 \cdot 6 = 2^4 \cdot 3 \cdot 5 \\1 \cdot 2 \cdot 3 \cdot 5 \cdot 6 = 2^2 \cdot 3^2 \cdot 5 \\1 \cdot 2 \cdot 3 \cdot 4 \cdot 6 = 2^4 \cdot 3^2 \\1 \cdot 2 \cdot 3 \cdot 4 \cdot 5 = 2^3 \cdot 3 \cdot 5\end{align*}


 The largest number that is certain to divide product P is basically GCD of all the above 6 products which is $2^2 \cdot 3$.

 Hence $P=3\cdot2^2=\boxed{\mathrm{(B)}\ 12}$.

 


	\item (\textit{AMC 10 10B 2004}) In the expression $c\cdot a^b-d$, the values of $a$, $b$, $c$, and $d$ are $0$, $1$, $2$, and $3$, although not necessarily in that order. What is the maximum possible value of the result?


 $\mathrm{(A) \ } 5 \qquad \mathrm{(B) \ } 6\qquad \mathrm{(C) \ } 8 \qquad \mathrm{(D) \ } 9\qquad \mathrm{(E) \ } 10$


 
 \textbf{Answer: }D



\textbf{Solution 1}

If $a=0$ or $c=0$, the expression evaluates to $-d<0$. 
If $b=0$, the expression evaluates to $c-d\leq 2$. 
Case $d=0$ remains.In that case, we want to maximize $c\cdot a^b$ where $\{a,b,c\}=\{1,2,3\}$. Trying out the six possibilities we get that the greatest is $(a,b,c)=(3,2,1)$, where $c\cdot a^b=1\cdot 3^2=\boxed{\mathrm{(D)}\ 9}$.

 



\textbf{Solution 2}

We first minimize $d$ as $0$, since is it subtracted and is independent from the rest. Then we are left with $a$, $b$, $c$ is from the set ${1,2,3}$. Then after some logical tries we can get $1\cdot3^2-0=\boxed{\mathrm{(D)}\ 9}$ is the largest.

 


	\item (\textit{AMC 10 10B 2004}) Which of the following numbers is a perfect square?


 $\noindent \textbf{(A)}\ 98! \cdot 99! \qquad \textbf{(B)}\ 98! \cdot 100! \qquad \textbf{(C)}\ 99! \cdot 100! \qquad \textbf{(D)}\ 99! \cdot 101! \qquad \textbf{(E)}\ 100! \cdot 101!$


 
 \textbf{Answer: }C



\textbf{Solution 1}

Using the fact that $n! = n\cdot (n-1)!$, we can write:

 
 \begin{align} A&=98! \cdot (99\cdot 98!) = 99 \cdot (98!)^2 = 11\cdot3^2\cdot(98!)^2 \\ B&=100 \cdot 99 \cdot (98!)^2 = 11\cdot10^2\cdot3^2\cdot( 98!)^2 \\ C&=100\cdot (99!)^2 = 10^2\cdot (99!)^2\\ D&=101\cdot 100\cdot (99!)^2 = 101 \cdot 10^2 \cdot (99!)^2\\ E& =101\cdot (100!)^2 \end{align}

 We see that $\boxed{\mathrm{(C) \ } 99! \cdot 100!}$ is a square, and because $11$, and $101$ are primes, none of the other four choices are squares.

 


 



\textbf{Solution 2}

Notice that $99! \cdot 99!$ is a perfect square (obviously). Also, $100 = 10^2$. Thus, 
 \[(99!)^2 \cdot 10^2 = 99! \cdot 100!\] is a perfect square, so the answer is $\boxed{\mathrm{(C) \ } 99! \cdot 100!}$.

 


	\item (\textit{AMC 10 10B 2004}) On a trip from the United States to Canada, Isabella took $d$ U.S. dollars. At the border she exchanged them all, receiving $10$ Canadian dollars for every $7$ U.S. dollars. After spending $60$ Canadian dollars, she had $d$ Canadian dollars left. What is the sum of the digits of $d$?


 $\mathrm{(A) \ } 5 \qquad \mathrm{(B) \ } 6\qquad \mathrm{(C) \ } 7\qquad \mathrm{(D) \ } 8\qquad \mathrm{(E) \ } 9$


 
 \textbf{Answer: }A



\textbf{Solution 1}

Isabella had $60+d$ Canadian dollars. Setting up an equation we get $d=\frac{7}{10}\cdot(60+d)$, which solves to $d=140$, and the sum of digits of $d$ is $\boxed{\mathrm{(A)}\ 5}$.

 



\textbf{Solution 2}

Each time Isabella exchanges $7$ U.S. dollars, she gets $7$ Canadian dollars and $3$ Canadian dollars extra. Isabella received a total of $60$ Canadian dollars extra, therefore she exchanged $7$ U.S. dollars $\frac{60}{3}=20$ times. Thus $d=7\cdot20=140$, and the sum of the digits is $\boxed{\mathrm{(A)}\ 5}$.

 



\textbf{Solution 3}

Isabella started with $d$ dollars and ended with $d$ dollars. Therefore,

 $\frac{d}{7}\cdot10-60=d$.

 Solving for $d$, we obtain $d = 140$. The sum of the digits of $d$ is $\boxed{\mathrm{(A)}\ 5}$.

 


	\item (\textit{AMC 10 10B 2004}) Minneapolis-St. Paul International Airport is $8$ miles southwest of downtown St. Paul and $10$ miles southeast of downtown Minneapolis. Which of the following is closest to the number of miles between downtown St. Paul and downtown Minneapolis?


 $\mathrm{(A) \ } 13 \qquad \mathrm{(B) \ } 14\qquad \mathrm{(C) \ } 15\qquad \mathrm{(D) \ } 16\qquad \mathrm{(E) \ } 17$


 
 \textbf{Answer: }A



\textbf{Solution}

The directions "southwest" and "southeast" are orthogonal. Thus the described situation is a right triangle with legs $8$ miles and $10$ miles long. The hypotenuse length is $\sqrt{8^2 + 10^2}\approx12.8$, and thus the answer is $\boxed{\mathrm{(A)}\ 13}$.

 Without a calculator one can note that $8^2+10^2=164<169=13^2\Rightarrow\boxed{\mathrm{(A)}\ 13}$.

 


	\item (\textit{AMC 10 10B 2004}) A square has sides of length $10$, and a circle centered at one of its vertices has radius $10$. What is the area of the union of the regions enclosed by the square and the circle?


 $\mathrm{(A) \ } 200+25\pi \qquad \mathrm{(B) \ } 100+75\pi \qquad \mathrm{(C) \ } 75+100\pi \qquad \mathrm{(D) \ } 100+100\pi \qquad \mathrm{(E) \ } 100+125\pi$


 
 \textbf{Answer: }B



\textbf{Solution}

The area of the circle is $S_{\bigcirc}=100\pi$; the area of the square is $S_{\square}=100$.

 Exactly $\frac{1}{4}$ of the circle lies inside the square. Thus the total area is $\dfrac34 S_{\bigcirc}+S_{\square}=\boxed{\mathrm{(B)\ }100+75\pi}$.

 
                    \begin{figure}[H]
                        \centering
                        \includegraphics[width=0.5\textwidth]{Solution/AMC_10_10B_2004_9_Solution_Figure_1.png}
                    \end{figure}
                    

 


	\item (\textit{AMC 10 10B 2004}) A grocer makes a display of cans in which the top row has one can and each lower row has two more cans than the row above it. If the display contains $100$ cans, how many rows does it contain?


 $\mathrm{(A) \ } 5 \qquad \mathrm{(B) \ } 8 \qquad \mathrm{(C) \ } 9 \qquad \mathrm{(D) \ } 10\qquad \mathrm{(E) \ } 11$


 
 \textbf{Answer: }D



\textbf{Solution}

The sum of the first $n$ odd numbers is $n^2$. As in our case $n^2=100$, we have $n=\boxed{\mathrm{(D)\ }10}$.

 


	\item (\textit{AMC 10 10B 2004}) Two eight-sided dice each have faces numbered $1$ through $8$. When the dice are rolled, each face has an equal probability of appearing on the top. What is the probability that the product of the two top numbers is greater than their sum?


 $\mathrm{(A) \ } \frac{1}{2} \qquad \mathrm{(B) \ } \frac{47}{64} \qquad \mathrm{(C) \ } \frac{3}{4} \qquad \mathrm{(D) \ } \frac{55}{64} \qquad \mathrm{(E) \ } \frac{7}{8}$


 
 \textbf{Answer: }C



\textbf{Solution 1}

We have $1\times n = n < 1 + n$, hence if at least one of the numbers is $1$, the sum is larger. There $15$ such possibilities.

 We have $2\times 2 = 2+2$.

 For $n>2$ we already have $2\times n = n + n > 2 + n$, hence all other cases are good.

 Out of the $8\times 8$ possible cases, we find that in $15+1=16$ the sum is greater than or equal to the product, hence in $64-16=48$ cases the sum is smaller, satisfying the condition. Therefore the answer is $\frac{48}{64}=\boxed{\mathrm{(C)\ }\frac{3}{4}}$.

 



\textbf{Solution 2}

Let the two rolls be $m$, and $n$. 

 From the restriction: $mn > m + n$

 $mn - m - n > 0$

 $mn - m - n + 1 > 1$

 $(m-1)(n-1) > 1$

 Since $m-1$ and $n-1$ are non-negative integers between $0$ and $7$, either  $(m-1)(n-1) = 0$, $(m-1)(n-1) = 1$, or $(m-1)(n-1) > 1$

 $(m-1)(n-1) = 0$ if and only if $m=1$ or $n=1$. 

 There are $8$ ordered pairs $(m,n)$ with $m=1$, $8$ ordered pairs with $n=1$, and $1$ ordered pair with $m=1$ and $n=1$. So, there are $8+8-1 = 15$ ordered pairs $(m,n)$ such that $(m-1)(n-1) = 0$. 

 $(m-1)(n-1) = 1$ if and only if $m-1=1$ and $n-1=1$ or equivalently $m=2$ and $n=2$. This gives $1$ ordered pair $(m,n) = (2,2)$. 

 So, there are a total of $15+1=16$ ordered pairs $(m,n)$ with $(m-1)(n-1) \le 1$. 

 Since there are a total of $8\cdot8 = 64$ ordered pairs $(m,n)$, there are $64-16 = 48$ ordered pairs $(m,n)$ with $(m-1)(n-1) > 1$. 

 Thus, the desired probability is $\frac{48}{64}=\boxed{\mathrm{(C)\ }\frac{3}{4}}$.

 


	\item (\textit{AMC 10 10B 2004}) An annulus is the region between two concentric circles. The concentric circles in the figure have radii $b$ and $c$, with $b>c$. Let $OX$ be a radius of the larger circle, let $XZ$ be tangent to the smaller circle at $Z$, and let $OY$ be the radius of the larger circle that contains $Z$. Let $a=XZ$, $d=YZ$, and $e=XY$. What is the area of the annulus? 


 
                    \begin{figure}[H]
                        \centering
                        \includegraphics[width=0.5\textwidth]{Problem/AMC_10_10B_2004_12_Problem_Figure_1.png}
                    \end{figure}
                    


 $\mathrm{(A) \ } \pi a^2 \qquad \mathrm{(B) \ } \pi b^2 \qquad \mathrm{(C) \ } \pi c^2 \qquad \mathrm{(D) \ } \pi d^2 \qquad \mathrm{(E) \ } \pi e^2$


 
 \textbf{Answer: }A



\textbf{Solution}

The area of the large circle is $\pi b^2$, the area of the small one is $\pi c^2$, hence the shaded area is $\pi(b^2-c^2)$.

 From the \href{https://artofproblemsolving.com/wiki/index.php?title=Pythagorean\_Theorem}{Pythagorean Theorem} for the right triangle $OXZ$ we have $a^2 + c^2 = b^2$, hence $b^2-c^2=a^2$ and thus the shaded area is $\boxed{\mathrm{(A)\ }\pi a^2}$.

 



\textbf{Solution 2}

Set $c=0,$ then the shaded area is just the area of a circle with radius $a,$ which is $\boxed{\mathrm{(A)\ }\pi a^2}$.

 


	\item (\textit{AMC 10 10B 2004}) In the United States, coins have the following thicknesses: penny, $1.55$ mm; nickel, $1.95$ mm; dime, $1.35$ mm; quarter, $1.75$ mm. If a stack of these coins is exactly $14$ mm high, how many coins are in the stack?


 $\mathrm{(A) \ } 7 \qquad \mathrm{(B) \ } 8 \qquad \mathrm{(C) \ } 9 \qquad \mathrm{(D) \ } 10 \qquad \mathrm{(E) \ } 11$


 
 \textbf{Answer: }B



\textbf{Solution 1}

All numbers in this solution will be in hundredths of a millimeter.

 The thinnest coin is the dime, with thickness $135$. A stack of $n$ dimes has height $135n$.

 The other three coin types have thicknesses $135+20$, $135+40$, and $135+60$. By replacing some of the dimes in our stack by other, thicker coins, we can clearly create exactly all heights in the set $\{135n, 135n+20, 135n+40, \dots, 195n\}$.

 If we take an odd $n$, then all the possible heights will be odd, and thus none of them will be $1400$. Hence $n$ is even.

 If $n<8$ the stack will be too low and if $n>10$ it will be too high. Thus we are left with cases $n=8$ and $n=10$.

 If $n=10$ the possible stack heights are $1350,1370,1390,\dots$, with the remaining ones exceeding $1400$.

 Therefore there are $\boxed{\mathrm{(B)\ }8}$ coins in the stack.

 Using the above observation we can easily construct such a stack. A stack of $8$ dimes would have height $8\cdot 135=1080$, thus we need to add $320$.This can be done for example by replacing five dimes by nickels (for $+60\cdot 5 = +300$), and one dime by a penny (for $+20$).

 



\textbf{Solution 2}

Let $p,n,d$, and $q$ be the number of pennies, nickels, dimes, and quarters used in the stack. 

 From the conditions above, we get the following equation:

 
 \[155p+195n+135d+175q=1400.\]

 Then we divide each side by five to get

 
 \[31p+39n+27d+35q=280.\]

 Writing both sides in terms of mod 4, we have $-p-n-d-q \equiv 0  \pmod 4$. 

 This means that the sum $p+n+d+q$ is divisible by 4. Therefore, the answer must be $\boxed{(B)\,\, 8}.$

 



\textbf{Solution 3}

We notice that the thickness of 4 quarters is 7 mm. 7 is half of 14, so we multiply the 4 quarters by two and get $\boxed{(B)\,\, 8}$.

 \textbf{Note} We can easily add up $1.55\text{\ mm}$ and $1.95\text{\ mm}$ to get $3.50\text{\ mm}$. We multiply that by $4$ to get $14\text{\ mm}$. Since this works and it requires 8 coins, the answer is clearly $\boxed{\mathrm{(B)\ }8}$.

 Similarly, we can simply take $8$ quarters to get $8\cdot 1.75=14$.

 


	\item (\textit{AMC 10 10B 2004}) A bag initially contains red marbles and blue marbles only, with more blue than red. Red marbles are added to the bag until only $\frac{1}{3}$ of the marbles in the bag are blue. Then yellow marbles are added to the bag until only $\frac{1}{5}$ of the marbles in the bag are blue. Finally, the number of blue marbles in the bag is doubled. What fraction of the marbles now in the bag are blue?


 $\mathrm{(A) \ } \frac{1}{5} \qquad \mathrm{(B) \ } \frac{1}{4} \qquad \mathrm{(C) \ } \frac{1}{3} \qquad \mathrm{(D) \ } \frac{2}{5} \qquad \mathrm{(E) \ } \frac{1}{2}$


 
 \textbf{Answer: }C



\textbf{Solution}

We can ignore most of the problem statement. The only important information is that immediately before the last step blue marbles formed $\frac{1}{5}$ of the marbles in the bag. This means that there were $x$ blue and $4x$ other marbles, for some $x$. When we double the number of blue marbles, there will be $2x$ blue and $4x$ other marbles, hence blue marbles now form $\boxed{\mathrm{(C)\ }\frac{1}{3}}$ of all marbles in the bag.

 



\textbf{Solution 2}

Assume that there are $2$ blue marbles and $1$ red marble. You can add $3$ red marbles to make the ratio of blue marbles to red marbles $\frac{1}{3}$. Then, to make the blue marbles $\frac{1}{5}$, add $4$ yellow marbles. Doubling the blue marbles makes $4$ blue marbles out of $12$ total marbles, or $\boxed{\mathrm{(C)\ }\frac{1}{3}}$ of all marbles in the bag.

 



\textbf{Solution 3}

If there are initially $B$ blue marbles in the bag, after red marbles are added, then the total number of marbles in the bag must be $3B$. Then after the yellow marbles are added, the number of marbles in the bag must be $5B$. Finally, adding $B$ blue marbles to the bag gives $2B$ blue marbles out of $6B$ total marbles. Thus $\boxed{1/3}$ of the marbles are blue.

 



\textbf{Solution 4}

Just before the number of blue marbles is doubled, the ratio of blue marbles to non-blue marbles is 1 to 4. Doubling the number of blue marbles makes the ratio 2 to 4, so $\dfrac{2}{2+4} = \boxed{1/3}$ of the marbles are blue.

 


	\item (\textit{AMC 10 10B 2004}) Patty has $20$ coins consisting of nickels and dimes. If her nickels were dimes and her dimes were nickels, she would have $70$ cents more. How much are her coins worth?


 $\textbf{(A)}\ \textdollar 1.15\qquad\textbf{(B)}\ \textdollar 1.20\qquad\textbf{(C)}\ \textdollar 1.25\qquad\textbf{(D)}\ \textdollar 1.30\qquad\textbf{(E)}\ \textdollar 1.35$


 
 \textbf{Answer: }A



\textbf{Solution 1}

She has $n$ nickels and $d=20-n$ dimes. Their total cost is $5n+10d=5n+10(20-n)=200-5n$ cents. If the dimes were nickels and vice versa, she would have $10n+5d=10n+5(20-n)=100+5n$ cents. This value should be $70$ cents more than the previous one. We get $200-5n+70=100+5n$, which solves to $n=17$. Her coins are worth $200-5n=\boxed{\mathrm{(A)\ }\textdollar1.15}$.

 



\textbf{Solution 2}

Changing a nickel into a dime increases the sum by $5$ cents, and changing a dime into a nickel decreases it by the same amount. As the sum increased by $70$ cents, there are $70/5=14$ more nickels than dimes. As the total count is $20$, this means that there are $17$ nickels and $3$ dimes, which is equal to $\boxed{\mathrm{(A)\ }\textdollar1.15}$.

 


	\item (\textit{AMC 10 10B 2004}) Three circles of radius $1$ are externally tangent to each other and internally tangent to a larger circle. What is the radius of the large circle?


 $\mathrm{(A) \ } \frac{2 + \sqrt{6}}{3} \qquad \mathrm{(B) \ } 2 \qquad \mathrm{(C) \ } \frac{2 + 3\sqrt{2}}{2} \qquad \mathrm{(D) \ } \frac{3 + 2\sqrt{3}}{3} \qquad \mathrm{(E) \ } \frac{3 + \sqrt{3}}{2}$


 
 \textbf{Answer: }D



\textbf{Solution 1}

The situation is shown in the picture below. The radius we seek is $SD = AD + AS$. Clearly $AD=1$. The point $S$ is the center of the equilateral triangle $ABC$, thus $AS$ is $2/3$ of the altitude of this triangle. We get that $AS = \frac23 \cdot \sqrt 3$. Therefore the radius we seek is$1 + \frac23 \cdot \sqrt 3 = \boxed{\mathrm{(D)\ }\frac{3+2\sqrt{3}}3}$.

 

                    \begin{figure}[H]
                        \centering
                        \includegraphics[width=0.5\textwidth]{Solution/AMC_10_10B_2004_16_Solution_Figure_1.png}
                    \end{figure}
                    

 



\textbf{Solution 2}

Using \href{https://artofproblemsolving.com/wiki/index.php?title=Descartes\%27\_Circle\_Formula}{Descartes' Circle Formula}, we can assign curvatures to all the circles: $1$, $1$, $1$, and $-\frac{1}{r}$ (b/c the bigger circle is externally tangent to all the other circles, the radius of the bigger circle is negative). Then, we can solve:

 $2(1^2+1^2+1^2+(-\frac{1}{r})^2) = (1+1+1-\frac{1}{r})^2$

 $r = \boxed{\mathrm{(D)\ }\frac{3+2\sqrt{3}}3}$

 


	\item (\textit{AMC 10 10B 2004}) The two digits in Jack's age are the same as the digits in Bill's age, but in reverse order. In five years Jack will be twice as old as Bill will be then. What is the difference in their current ages?


 $\mathrm{(A) \ } 9 \qquad \mathrm{(B) \ } 18 \qquad \mathrm{(C) \ } 27 \qquad \mathrm{(D) \ } 36\qquad \mathrm{(E) \ } 45$


 
 \textbf{Answer: }B



\textbf{Solution 1}

If Jack's current age is $\overline{ab}=10a+b$, then Bill's current age is $\overline{ba}=10b+a$.

 In five years, Jack's age will be $10a+b+5$ and Bill's age will be $10b+a+5$. 

 We are given that $10a+b+5=2(10b+a+5)$. 

 
Thus $8a=19b+5 \Rightarrow a=\dfrac{19b+5}{8}$. 

 For $b=1$ we get $a=3$. For $b=2$ and $b=3$ the value $\frac{19b+5}8$ is not an integer, and for $b\geq 4$, $a$ is more than $9$. Thus the only solution is $(a,b)=(3,1)$, and the difference in ages is $31-13=\boxed{\mathrm{(B)\ }18}$.

 



\textbf{Solution 2}

Age difference does not change in time. Thus in five years Bill's age will be equal to their age difference.

 The age difference is $(10a+b)-(10b+a)=9(a-b)$, hence it is a multiple of $9$. Thus Bill's current age modulo $9$ must be $4$.

 Thus Bill's age is in the set $\{13,22,31,40,49,58,67,76,85,94\}$.

 As Jack is older, we only need to consider the cases where the tens digit of Bill's age is smaller than the ones digit. This leaves us with the options $\{13,49,58,67\}$. 

 Checking each of them, we see that only $13$ works, and gives the solution $31-13=\boxed{\mathrm{(B)\ }18}$.

 


	\item (\textit{AMC 10 10B 2004}) In the right triangle $\triangle ACE$, we have $AC=12$, $CE=16$, and $EA=20$. Points $B$, $D$, and $F$ are located on $AC$, $CE$, and $EA$, respectively, so that $AB=3$, $CD=4$, and $EF=5$. What is the ratio of the area of $\triangle DBF$ to that of $\triangle ACE$?


 
                    \begin{figure}[H]
                        \centering
                        \includegraphics[width=0.5\textwidth]{Problem/AMC_10_10B_2004_18_Problem_Figure_1.png}
                    \end{figure}
                    


 $\mathrm{(A) \ } \frac{1}{4} \qquad \mathrm{(B) \ } \frac{9}{25} \qquad \mathrm{(C) \ } \frac{3}{8} \qquad \mathrm{(D) \ } \frac{11}{25} \qquad \mathrm{(E) \ } \frac{7}{16}$


 
 \textbf{Answer: }E



\textbf{Solution 1 (Trigonometry)}

Let $x = [DBF]$. Because $\triangle ACE$ is divided into four triangles, $[ACE] = [BCD] + [ABF] + [DEF] + x$.

 Because of $SAS$ triangle area, $\frac12 \cdot 12 \cdot 16 = \frac12 \cdot 9 \cdot 4 + \frac12 \cdot 3 \cdot 15 \cdot \sin(\angle A) + \frac12 \cdot 5 \cdot 12 \cdot \sin(\angle E) + x$.

 $\sin(\angle A) = \frac{16}{20}$ and $\sin(\angle E) = \frac{12}{20}$, so $96 = 18 + 18 + 18 + x$.

 $x = 42$, so $\frac{[DBF]}{[ACE]} = \frac{42}{96} = \boxed{\textbf{(E)}\frac 7{16}}$.

 



\textbf{Solution 2}

First of all, note that $\frac{AB}{AC} = \frac{CD}{CE} = \frac{EF}{EA} = \frac 14$, and therefore$\frac{BC}{AC} = \frac{DE}{CE} = \frac{FA}{EA} = \frac 34$.

 Draw the height from $F$ onto $AB$ as in the picture below:

 
                    \begin{figure}[H]
                        \centering
                        \includegraphics[width=0.5\textwidth]{Solution/AMC_10_10B_2004_18_Solution_Figure_1.png}
                    \end{figure}
                    

 Now consider the area of $\triangle ABF$. Clearly the triangles $\triangle AFG$ and $\triangle AEC$ are similar, as they have all angles equal. Their ratio is $\frac {AF}{AE} = \frac 34$, hence $FG = \frac 34 \cdot CE$. Now the area $S_{ABF}$ of $\triangle ABF$ can be computed as $S_{ABF} = \frac 12 \cdot AB \cdot FG$ = $\frac 12 \cdot \left( \frac 14 \cdot AC \right) \cdot \left( \frac 34 \cdot EC \right) = \frac 14 \cdot \frac 34 \cdot S_{ACE}$.

 Similarly we can find that $S_{BCD} = S_{DEF} = \frac 3{16}\cdot S_{ACE}$ as well. 

 Hence $S_{BDF} = S_{ACE} - 3\cdot\left( \frac 3{16} \cdot S_{ACE} \right) = \frac 7{16} \cdot S_{ACE}$, and the answer is $\frac{S_{BDF}}{S_{ACE}} = \boxed{\frac 7{16}}$.

 



\textbf{Solution 3 (Coordinate Geometry)}

We will put triangle ACE on a xy-coordinate plane with C being the origin. The area of triangle ACE is 96. To find the area of triangle DBF, let D be (4, 0), let B be (0, 9), and let F be (12, 3). You can then use the shoelace theorem to find the area of DBF, which is 42. $\frac {42}{96} = \boxed{\frac 7{16}}$

 


 



\textbf{Solution 4}

You can also place a point $X$ on $CE$ such that $CX$ is $12$, creating trapezoid $CBFX$. Then, you can find the area of the trapezoid, subtract the area of the two right triangles $DFX$ and $BCD$, divide by the area of $ABC$, and get the ratio of $7/16$.

 



\textbf{Solution 5}

It is well known that for when two triangles share an angle, the two sides around the shared angle is proportional to the areas of each of the two triangles. 

 We can find all the ratios of the triangles except for $\triangle{BDF}$ and then subtract from $1.$

 In this case, we have $\triangle{ABF}$ sharing $\angle{A}$ with $\triangle{ACE}$.

 Therefore, we have $\frac{[ABF]}{[ACE]}=\frac{AB}{AC} \cdot \frac{AF}{AE} = \frac{3}{12} \cdot \frac{15}{20} = \frac{3}{16}.$

 Also note that $\triangle{EFD}$ shares $\angle{E}$ with $\triangle{EAC}$.

 Therefore, we have $\frac{[EFD]}{[EAC]}=\frac{ED}{EC} \cdot \frac{EF}{EA} = \frac{12}{16} \cdot \frac{5}{20} = \frac{3}{16}.$

 Lastly, note that $\triangle{CDB}$ shares $\angle{C}$ with $\triangle{CAE}$.

 Therefore, we have $\frac{[CDB]}{[CAE]}=\frac{CD}{CE} \cdot \frac{CB}{CA} = \frac{9}{12} \cdot \frac{4}{16} = \frac{3}{16}.$

 Thus, the ratio of $\triangle{DBF}$ to $\triangle{ACE}$ is $1- \left( \frac{3}{16}+\frac{3}{16}+\frac{3}{16} \right)=1-\frac{9}{16}=\boxed{\textbf{(E)} \frac {7}{16}}$

 



\textbf{Solution 6 (Wooga Looga Theorem)}

We know that $\frac{CB}{BA}=\frac{AF}{FE}=\frac{ED}{DC}=3$, so by the \href{https://artofproblemsolving.com/wiki/index.php?title=The\_Devil\%27s\_Triangle}{The Devil's Triangle} we have $\frac{[DBF]}{[ACE]}=\frac{3^2-3+1}{(3+1)^2}=\boxed{\frac7{16}}$.

 


	\item (\textit{AMC 10 10B 2004}) In the sequence $2001$, $2002$, $2003$, $\ldots$ , each term after the third is found by subtracting the previous term from the sum of the two terms that precede that term. For example, the fourth term is $2001 + 2002 - 2003 = 2000$. What is the
$2004^\textrm{th}$ term in this sequence?


 $\mathrm{(A) \ } -2004 \qquad \mathrm{(B) \ } -2 \qquad \mathrm{(C) \ } 0 \qquad \mathrm{(D) \ } 4003 \qquad \mathrm{(E) \ } 6007$


 
 \textbf{Answer: }C



\textbf{Solution 1}

We already know that $a_1=2001$, $a_2=2002$, $a_3=2003$, and $a_4=2000$. Let's compute the next few terms to get the idea how the sequence behaves. We get $a_5 = 2002+2003-2000 = 2005$, $a_6=2003+2000-2005=1998$, $a_7=2000+2005-1998=2007$, and so on.

 We can now discover the following pattern: $a_{2k+1} = 2001+2k$ and $a_{2k}=2004-2k$. This is easily proved by induction. It follows that $a_{2004}=a_{2\cdot 1002} = 2004 - 2\cdot 1002 = \boxed{0}$.

 



\textbf{Solution 1 but in a more not smart way}

Subtract 2000 from each of the terms so the sequence turns into 1, 2, 3, 0, 5, -2, 7, -4, 9, -6... (1+2-3=0, 2+3-0=5, 3+0-5=-2, 0+5-(-2)=7, 5+(-2)-7=-4, etc.). Quickly notice that after taking away 1 and 2, the 2nd term is 0, the 4th term is -2, the 6th term is -4, the 8th term is -6, etc. Thus, in the sequence without 1 and 2, the (2n)th term has a value of -2(n-1). Therefore, to find the 2004th term in the original sequence, we take away the first 2 values to form the new sequence so the value of the 2004th term in the original sequence is the 2002nd term in the new sequence. 2002 is 2 times 1001. Thus, the value is -2(1001-1)= -2000. BUT, REMEMBER THAT AT THE START WE TOOK AWAY 2000 FROM EACH TERM. ADD -2000 TO 2000 TO GET $\boxed{0}$. LOLOL

 



\textbf{Solution 2}

Note that the recurrence $a_n+a_{n+1}-a_{n+2}~=~a_{n+3}$ can be rewritten as $a_n+a_{n+1} ~=~ a_{n+2}+a_{n+3}$. 

 Hence we get that $a_1+a_2 ~=~ a_3+a_4 ~=~ a_5+a_6 ~= \cdots$ and also$a_2+a_3 ~=~ a_4+a_5 ~=~ a_6+a_7 ~= \cdots$

 From the values given in the problem statement we see that $a_3=a_1+2$. 

 From $a_1+a_2 = a_3+a_4$ we get that $a_4=a_2-2$.

 From $a_2+a_3 = a_4+a_5$ we get that $a_5=a_3+2$.

 Following this pattern, we get $a_{2004} = a_{2002} - 2 = a_{2000} - 4 = \cdots = a_2 - 2002 = \boxed{0}$.

 



\textbf{Solution 3}

Our recurrence is $a_n+a_{n+1}-a_{n+2}~=~a_{n+3}$, so we get $r^3+r^2-r-1 = 1$, so $(r-1)(r+1)^2 = 1$, so our formula for the recurrence is $a_n = A + (B + Cn)(-1)^n$.

 Substituting our starting values gives us $a_n = 2002 + (2 - n)(-1)^n$.

 So, $a_{2004} = 2002 - 2002 = 0.$

 


	\item (\textit{AMC 10 10B 2004}) In $\triangle ABC$ points $D$ and $E$ lie on $BC$ and $AC$, respectively. If $AD$ and $BE$ intersect at $T$ so that $\frac{AT}{DT}=3$ and $\frac{BT}{ET}=4$, what is $\frac{CD}{BD}$?


 
                    \begin{figure}[H]
                        \centering
                        \includegraphics[width=0.5\textwidth]{Problem/AMC_10_10B_2004_20_Problem_Figure_1.png}
                    \end{figure}
                    


 $\mathrm{(A) \ } \frac{1}{8} \qquad \mathrm{(B) \ } \frac{2}{9} \qquad \mathrm{(C) \ } \frac{3}{10} \qquad \mathrm{(D) \ } \frac{4}{11} \qquad \mathrm{(E) \ } \frac{5}{12}$


 
 \textbf{Answer: }D



\textbf{Solution 1 (Triangle Areas)}

We use the square bracket notation $[\cdot]$ to denote area.

 WLOG, we can assume $[\triangle BTD] = 1$. Then $[\triangle BTA] = 3$, and $[\triangle ATE] = 3/4$. We have $CD/BD = [\triangle  ACD]/[\triangle ABD]$, so we need to find the area of quadrilateral $TDCE$.

 Draw the line segment $TC$ to form the two triangles $\triangle TDC$ and $\triangle TEC$. Let $x = [\triangle TDC]$, and $y = [\triangle TEC]$. By considering triangles $\triangle BTC$ and $\triangle ETC$, we obtain $(1+x)/4=y$, and by considering triangles $\triangle ATC$ and $\triangle DTC$, we obtain $\dfrac{\frac{3}{4} +y}{3}=x$. Solving, we get $x=4/11$, $y=15/44$, so the area of quadrilateral $TDCE$ is $x+y=31/44$.

 Therefore $\frac{CD}{BD}=\frac{\frac{3}{4}+\frac{31}{44}}{3+1}=\boxed{\textbf{(D)} \frac{4}{11}}$

 



\textbf{Solution 2 (Triangle Areas, Alternate Approach)}

We observe that $\frac{BC}{CD} = \frac{[BAE]}{[DAE]}$.  The proof is that if $B_H$ and $D_H$ are the feet of the altitudes from $B$ and $D$, respectively, to $AC$, then both sides of that equation are equal to $\frac{BB_H}{DD_H}$.

 From there, we can easily finish:
 \[\frac{BC}{CD} = \frac{[BAE]}{[DAE]} = \frac{[BAT]\cdot\frac{5}{4}}{[EAT]\cdot \frac{4}{3}} = 4\cdot\frac{5}{4}\cdot\frac{3}{4} = \frac{15}{4}\]and thus $\frac{CD}{BD} = \boxed{\textbf{(D)}\frac 4{11}}$.

 



\textbf{Solution 3 (Mass points)}

The presence of only ratios in the problem essentially cries out for mass points.

 As per the problem, we assign a mass of $1$ to point $A$, and a mass of $3$ to $D$. Then, to balance $A$ and $D$ on $T$, $T$ has a mass of $4$.

 Now, were we to assign a mass of $1$ to $B$ and a mass of $4$ to $E$, we'd have $5T$. Scaling this down by $4/5$ (to get $4T$, which puts $B$ and $E$ in terms of the masses of $A$ and $D$), we assign a mass of $\frac{4}{5}$ to $B$ and a mass of $\frac{16}{5}$ to $E$.

 Now, to balance $A$ and $C$ on $E$, we must give $C$ a mass of $\frac{16}{5}-1=\frac{11}{5}$.

 Finally, the ratio of $CD$ to $BD$ is given by the ratio of the mass of $B$ to the mass of $C$, which is $\frac{4}{5}\cdot\frac{5}{11}=\boxed{\textbf{(D)}\ \frac{4}{11}}$.

 



\textbf{Solution 4 (Coordinates)}

Affine transformations preserve ratios of distances, and for any pair of triangles, there is an affine transformation that maps the first one onto the second one. This is why the answer is the same for any $\triangle ABC$, and we just need to compute it for any single triangle.

 We can choose the points $A=(-3,0)$, $B=(0,4)$, and $D=(1,0)$. This way we will have $T=(0,0)$, and $E=(0,-1)$. The situation is shown in the picture below:

 
                    \begin{figure}[H]
                        \centering
                        \includegraphics[width=0.5\textwidth]{Solution/AMC_10_10B_2004_20_Solution_Figure_1.png}
                    \end{figure}
                    

 The point $C$ is the intersection of the lines $BD$ and $AE$. The points on the first line have the form $(t,4-4t)$, the points on the second line have the form $(t,-1-t/3)$. Solving for $t$ we get $t=15/11$, hence $C=(15/11,-16/11)$.

 The ratio $CD/BD$ can now be computed simply by observing the $x$ coordinates of $B$, $C$, and $D$:

 
 \[\frac{CD}{BD} = \frac{15/11 - 1}{1 - 0} = \boxed{\textbf{(D)}\frac 4{11}}\]

 



\textbf{Solution 5 (Menelaus)}

By Menelaus on triangle $BDT$, we have that 
 \begin{align*}1&=\dfrac{DA}{TA}\cdot\dfrac{TE}{BE}\cdot\dfrac{BC}{DC} \\ &= \dfrac43 \cdot \dfrac 15 \cdot \dfrac {BC}{DC},\end{align*}
giving $\dfrac{BC}{DC} = \dfrac{15}{4}$. Therefore, $\dfrac{CD}{BD} = \boxed{\textbf{(D)}\dfrac{4}{11}}$.

 



\textbf{Solution 6 (vector)}

Since $\frac{BT}{ET} = 4$, we have 
 \[\overrightarrow{AT} = \frac{1}{5}\overrightarrow{AB} + \frac{4}{5}\overrightarrow{AE}\]Furthermore, since $\frac{AT}{DT} = 3$, it follows that 
 \[\overrightarrow{AD} = \frac{4}{3}\overrightarrow{AT}\]Therefore, 
 \[\overrightarrow{AD} = \frac{4}{15}\overrightarrow{AB} + \frac{16}{15}\overrightarrow{AE}\]

 Let $\overrightarrow{AC} = t\overrightarrow{AE}$. Then 
 \[\overrightarrow{AD} = \frac{4}{15}\overrightarrow{AB} + \frac{16}{15t}\overrightarrow{AC}\]Since points $D$, $B$, and $C$ are collinear, 
 \[\frac{4}{15} + \frac{16}{15t} = 1 \implies t = \frac{16}{11}\]Thus, 
 \[\overrightarrow{AD} = \frac{4}{15}\overrightarrow{AB} + \frac{11}{15}\overrightarrow{AC}\]which yields 
 \[\frac{CD}{BD} = \frac{4}{11}\]

 


	\item (\textit{AMC 10 10B 2004}) Let $1$; $4$; $\ldots$ and $9$; $16$; $\ldots$ be two arithmetic progressions. The set $S$ is the union of the first $2004$ terms of each sequence. How many distinct numbers are in $S$?


 $\mathrm{(A) \ } 3722 \qquad \mathrm{(B) \ } 3732 \qquad \mathrm{(C) \ } 3914 \qquad \mathrm{(D) \ } 3924 \qquad \mathrm{(E) \ } 4007$


 
 \textbf{Answer: }A



\textbf{Solution 1}

The terms in the first sequence are defined by $n \equiv 1 \pmod3$, while the terms in the second sequence are defined by $n \equiv 2 \pmod 7.$ We seek to find the solutions to this system of modular congruences.

 Letting $n = 3a + 1 = 7b + 2$ for nonnegative integers $a,b$ we solve the congruence \begin{align*} 3a+1 &\equiv 7b+2 \pmod 3, \\ 1 &\equiv b + 2 \pmod 3, \\ b &\equiv 2 \pmod 3.\end{align*}


 For nonnegative integer $c,$ we have that $b = 3c+2$. Substituting back into our original equation, we have $n = 21c + 16,$ or $n \equiv 16 \pmod {21}.$

 Now we have to find the largest term in the smaller sequence (first sequence), which is $2003 \cdot 3 + 1 = 6010.$ Dividing $\frac{6010}{21}$ gives that there are $286$ terms in common between the sequences.

 By PIE, the answer is just $2004 + 2004 - 286 = \boxed{(A) 3722}.$

 



\textbf{Solution 2}

The two sets of terms are $A=\{ 3k+1 : 0\leq k < 2004 \}$ and $B=\{ 7l+9 : 0\leq l<2004\}$. 

 Now $S=A\cup B$. We can compute $|S|=|A\cup B|=|A|+|B|-|A\cap B|=4008-|A\cap B|$. We will now find $|A\cap B|$. 

 Consider the numbers in $B$. We want to find out how many of them lie in $A$. In other words, we need to find out the number of valid values of $l$ for which $7l+9\in A$.

 The fact "$7l+9\in A$" can be rewritten as "$1\leq 7l+9 \leq 3\cdot 2003 + 1$, and $7l+9\equiv 1\pmod 3$".

 The first condition gives $0\leq l\leq 857$, the second one gives $l\equiv 1\pmod 3$.

 Thus the good values of $l$ are $\{1,4,7,\dots,856\}$, and their count is $858/3 = 286$.

 Therefore $|A\cap B|=286$, and thus $|S|=4008-|A\cap B|=\boxed{(A) 3722}$.

 



\textbf{Solution 3}

We can start by finding the first non-distinct term from both sequences. We find that that number is $16$. Now, to find every other non-distinct terms, we can just keep adding $21$. 

 We know that the last terms of both sequences are $1+3\cdot 2003$ and $9+7\cdot 2003$. Clearly, $1+3\cdot 2003$ is smaller and that is the last possible common term of both sequences. 

 Now, we can create the inequality $16+21k \leq 1+3\cdot 2003$. Using the inequality, we find that there are $286$ common terms. There are 4008 terms in total. $4008-286=\boxed{(A) 3722}$

 


	\item (\textit{AMC 10 10B 2004}) A triangle with sides of length $5, 12,$ and $13$ has both an inscribed and a circumscribed circle. What is the distance between the centers of those circles?


 $\mathrm{(A) \ } \frac{3\sqrt{5}}{2} \qquad \mathrm{(B) \ } \frac{7}{2} \qquad \mathrm{(C) \ } \sqrt{15} \qquad \mathrm{(D) \ } \frac{\sqrt{65}}{2} \qquad \mathrm{(E) \ } \frac{9}{2}$


 
 \textbf{Answer: }D



\textbf{Solution 1}


                    \begin{figure}[H]
                        \centering
                        \includegraphics[width=0.5\textwidth]{Solution/AMC_10_10B_2004_22_Solution_Figure_1.png}
                    \end{figure}
                    This is a right triangle. Pick a coordinate system so that the right angle is at $(0,0)$ and the other two vertices are at $(12,0)$ and $(0,5)$.

 As this is a right triangle, the center of the circumcircle is in the middle of the hypotenuse, at $(6,2.5)$.

 The radius $r$ of the inscribed circle can be computed using the well-known identity $\frac{rP}2=S$, where $S$ is the area of the triangle and $P$ its perimeter. In our case, $S=\frac{5\cdot 12}{2}=30$ and $P=5+12+13=30$. Thus, $r=2$. As the inscribed circle touches both legs, its center must be at $(r,r)=(2,2)$.

 The distance of these two points is then $\sqrt{ (6-2)^2 + (2.5-2)^2 } = \sqrt{16.25} = \sqrt{\frac{65}4} = \boxed{\frac{\sqrt{65}}2}$.

 



\textbf{Solution 2}

We directly apply Euler’s Theorem, which states that if the circumcenter is $O$ and the incenter $I$, and the inradius is $r$ and the circumradius is $R$, then 
 \[OI^2=R(R-2r)\]

 We can see that this is a right triangle, and hence has area $30$. We then find the inradius with the formula $A=rs$, where $s$ denotes semiperimeter. We easily see that $s=15$, so $r=2$. 

 We now find the circumradius with the formula $A=\frac{abc}{4R}$. Solving for $R$ gives $R=\frac{13}{2}$. Additionally, we may notice that the side lengths $(5, 12, 13)$ are in a Pythagorean triple, and therefore the triangle is right for a circumradius of $\frac{13}{2}$. 

 Substituting all of this back into our formula gives:
 \[OI^2= \frac{65}{4}\]So, $OI=\frac{\sqrt{65}}{2}\implies \boxed{D}$

 



\textbf{Solution 3}


                    \begin{figure}[H]
                        \centering
                        \includegraphics[width=0.5\textwidth]{Solution/AMC_10_10B_2004_22_Solution_Figure_2.png}
                    \end{figure}
                    Construct $\triangle{ABC}$ such that $AB=5$, $AC=12$, and $BC=13$. Since this is a pythagorean triple, $\angle{A}=90$. By a property of circumcircles and right triangles, the circumcenter, $G$, lies on the midpoint of $\overline{BC}$, so $BG=\frac{13}{2}$. Turning to the incircle, we find that the inradius is $2$, using the formula $A=rs$, where $A$ is the area of the triangle, $r$ is the inradius, and $s$ is the semiperimeter. We then denote the incenter $I$, along with the points of tangency $D$, $E$, and $F$. Because $\angle{IDA}=\angle{IEA}=90$ by a property of tangency, $\angle{EID}=90$, and so $IDAE$ is a square. Then, since $IE=2$, $AD=2$. As $AB=5$, $BD=3$, and because $\triangle{BID}\cong\triangle{BIF}$ by HL, $BD=BF=3$. Therefore, $FG=\frac{7}{2}$. Because $IF=2$, pythagorean theorem gives $IG=\boxed{\frac{\sqrt{65}}{2}}$

 


 



\textbf{Solution 4 (Coordinate Bash)}

Like Solution 1, this is a right triangle, so use a coordinate system and put the right angle at $(0,0)$ and the other two vertices at $(12,0)$ and $(0,5)$. Call point A $(0,5)$, B $(0,0)$ , and C $(12,0)$. Then you need to find the center of the circumscribed and inscribed circles, or circumcenter and incenter of the triangle. Find the center of the circumcenter using the other methods above, which will be at $(6,2.5)$. In order to find the incenter, you will use the fact that it is the point where the three bisectors of the interior angles of the triangle intersect. From point B, the bisector is 45 degrees, so the slope is $y=x$. From point A, draw the angle bisector through the line $(0,0)$ and $(12,0)$ and call it point D. Now you need to find $BD$. Replacing $BD$ with $x$ and $DC$ with $y$ and using the angle bisector theorem $\frac{x}{y} = \frac{AB}{AC} = \frac{5}{13}$, so $y=\frac{13}{5}x$. Substituting $x+y=12$, we get $x+\frac{13}{5}x = \frac{18}{5}x = 12$, so $x = \frac{60}{18} = \frac{10}{3}$. With this information, point D is $(\frac{10}{3}, 0)$, and the slope of $BD$ is $\frac{rise}{run} = \frac{-5}{\frac{10}{3}} = -\frac{3}{2}$. This means the equation for $BD$ is $y=-\frac{3}{2}x+5$. Find the intersection point of $y=\frac{-3}{2}x+5$ and $y=x$ with substitution, which will be $(2,2)$. Using the distance formula, find the distance between $(2,2)$ and $(6,2.5)$ which will be $\sqrt{16.25}$, and with some manipulation, $\frac{\sqrt{16.25}\cdot\sqrt{4}}{\sqrt{4}} = \frac{\sqrt{65}}{2}\implies \boxed{D}$.

 
                \begin{itemize}
                    \item Someone please draw a diagram and try to make it cleaner to read.
                \end{itemize}


                


	\item (\textit{AMC 10 10B 2004}) Each face of a cube is painted either red or blue, each with probability $1/2$. The color of each face is determined independently. What is the probability that the painted cube can be placed on a horizontal surface so that the four vertical faces are all the same color?


 $\mathrm{(A) \ } \frac{1}{4} \qquad \mathrm{(B) \ } \frac{5}{16} \qquad \mathrm{(C) \ } \frac{3}{8} \qquad \mathrm{(D) \ } \frac{7}{16} \qquad \mathrm{(E) \ } \frac{1}{2}$


 
 \textbf{Answer: }B



\textbf{Solution 1}

Label the six sides of the cube by numbers $1$ to $6$ as on a classic dice. Then the "four vertical faces" can be: $\{1,2,5,6\}$, $\{1,3,4,6\}$, or $\{2,3,4,5\}$.

 Let $A$ be the set of colorings where $1,2,5,6$ are all of the same color, similarly let $B$ and $C$ be the sets of good colorings for the other two sets of faces.

 There are $2^6=64$ possible colorings, and there are $|A\cup B\cup C|$ good colorings. Thus the result is $\frac{|A\cup B\cup C|}{64}$. We need to compute $|A\cup B\cup C|$. 

 Using the \href{https://artofproblemsolving.com/wiki/index.php?title=Principle\_of\_Inclusion-Exclusion}{Principle of Inclusion-Exclusion} we can write
 \[|A\cup B\cup C| = |A|+|B|+|C| - |A\cap B| - |A\cap C| - |B\cap C| + |A\cap B\cap C|\]

 Clearly $|A|=|B|=|C|=2^3=8$, as we have two possibilities for the common color of the four vertical faces, and two possibilities for each of the horizontal faces.

 What is $A\cap B$? The faces $1,2,5,6$ must have the same color, and at the same time faces $1,3,4,6$ must have the same color. It turns out that $A\cap B=A\cap C=B\cap C= A\cap B\cap C =$ the set containing just the two cubes where all six faces have the same color.

 Therefore$|A\cup B\cup C| = 8+8+8-2-2-2+2 = 20$, and the result is $\frac{20}{64}=\boxed{\frac{5}{16}}$.

 



\textbf{Solution 2}

Suppose we break the situation into cases that contain four vertical faces of the same color:

 I. Two opposite sides of same color:There are 3 ways to choose the two sides, and then two colors possible, so $3 \cdot 2=6$.

 II. One face different from all the others:There are 6 ways to choose this face, and 2 colors, so $6 \cdot 2=12$.

 III. All faces are the same:There are 2 colors, and so two ways for all faces to be the same. 

 Adding them up, we have a total of $20$ ways to have four vertical faces the same color. There are $2^6$ ways to color the cube, so the answer is $\frac{20}{64}=\boxed{\frac{5}{16}}$.

 


	\item (\textit{AMC 10 10B 2004}) In $\bigtriangleup ABC$ we have $AB = 7$, $AC = 8$, and $BC = 9$. Point $D$ is on the circumscribed circle of the triangle so that $AD$ bisects $\angle BAC$. What is the value of $\frac{AD}{CD}$?


 $\mathrm{(A) \ } \frac{9}{8} \qquad \mathrm{(B) \ } \frac{5}{3} \qquad \mathrm{(C) \ } 2 \qquad \mathrm{(D) \ } \frac{17}{7} \qquad \mathrm{(E) \ } \frac{5}{2}$


 
 \textbf{Answer: }B



\textbf{Solution 1 (Ptolemy's Theorem)}

Set $\overline{BD}$'s length as $x$. $\overline{CD}$'s length must also be $x$ since $\angle BAD$ and $\angle DAC$ intercept arcs of equal length (because $\angle BAD=\angle DAC$). Using \href{https://artofproblemsolving.com/wiki/index.php?title=Ptolemy\%27s\_Theorem}{Ptolemy's Theorem}, $7x+8x=9(AD)$. The ratio is $\frac{5}{3}\implies\boxed{\text{(B)}}$

 



\textbf{Solution 2 (Similarity Proportion)}


                    \begin{figure}[H]
                        \centering
                        \includegraphics[width=0.5\textwidth]{Solution/AMC_10_10B_2004_24_Solution_Figure_1.png}
                    \end{figure}
                    Let $E = \overline{BC}\cap \overline{AD}$.  Observe that $\angle ABC \cong \angle ADC$ because they both subtend arc $\overarc{AC}.$ 

 Furthermore, $\angle BAE \cong \angle EAC$ because $\overline{AE}$ is an angle bisector, so $\triangle ABE \sim \triangle ADC$ by $\text{AA}$ similarity. Then $\dfrac{AD}{AB} = \dfrac{CD}{BE}$.  By the \href{https://artofproblemsolving.com/wiki/index.php?title=Angle\_Bisector\_Theorem}{Angle Bisector Theorem}, $\dfrac{7}{BE} = \dfrac{8}{CE}$, so $\dfrac{7}{BE} = \dfrac{8}{9-BE}$. This in turn gives $BE = \frac{21}{5}$.  Plugging this into the similarity proportion gives:  $\dfrac{AD}{7} = \dfrac{CD}{\frac{21}{5}} \implies \dfrac{AD}{CD} = {\dfrac{5}{3}} = \boxed{\text{(B)}}$.

 



\textbf{Solution 3 (Angle Bisector Theorem)}

Similar to solution 2, let $E$ be the intersection of diagonals $AD$ and $BC$. We know that $\overline{AD}$ bisects $\angle BAC$, so $\angle BAD = \angle CAD$. Additionally, $\angle BAD$ and $\angle BCD$ subtend the same arc, giving $\angle BAD = \angle BCD$. Similarly, $\angle CAD = \angle CBD$ and $\angle ABC = \angle ADC$.

 These angle relationships tell us that $\triangle ABE\sim \triangle ADC$ by AA Similarity, so $AD/CD = AB/BE$. By the angle bisector theorem, $AB/BE = AC/CE$. Hence,
 \[\frac{AB}{BE} = \frac{AC}{CE} = \frac{AB + AC}{BE + CE} = \frac{AB + AC}{BC} = \frac{7 + 8}{9} = \frac{15}{9} = \boxed{\frac{5}{3}}.\]

 --vaporwave

 


	\item (\textit{AMC 10 10B 2004}) A circle of radius $1$ is internally tangent to two circles of radius $2$ at points $A$ and $B$, where $AB$ is a diameter of the smaller circle. What is the area of the region, shaded in the picture, that is outside the smaller circle and inside each of the two larger circles? 


 $\mathrm{(A) \ } \frac{5}{3} \pi - 3\sqrt 2 \qquad  \mathrm{(B) \ } \frac{5}{3} \pi - 2\sqrt 3 \qquad  \mathrm{(C) \ } \frac{8}{3} \pi - 3\sqrt 3 \qquad  \mathrm{(D) \ } \frac{8}{3} \pi - 3\sqrt 2 \qquad  \mathrm{(E) \ } \frac{8}{3} \pi - 2\sqrt 3$


 


                    \begin{figure}[H]
                        \centering
                        \includegraphics[width=0.5\textwidth]{Problem/AMC_10_10B_2004_25_Problem_Figure_1.png}
                    \end{figure}
                    


 
 \textbf{Answer: }B



\textbf{Solution 1}

The area of the small circle is $\pi$. We can add it to the shaded region, compute the area of the new region, and then subtract the area of the small circle from the result.

 Let $C$ and $D$ be the intersections of the two large circles. Connect them to $A$ and $B$ to get the picture below:

 
                    \begin{figure}[H]
                        \centering
                        \includegraphics[width=0.5\textwidth]{Solution/AMC_10_10B_2004_25_Solution_Figure_1.png}
                    \end{figure}
                    

 We can see that the triangles $\triangle ABC$ and $\triangle ABD$ are both equilateral with side $2$.

 Take a look at the lower circle. The angle $ABC$ is $60^\circ$, thus sector $ABC$ is $1/6$ of the circle. The same is true for sector $ABD$ of the lower circle, and sectors $CAB$ and $BAD$ of the upper circle. 

 If we now sum the areas of these four sectors, we will almost get the area of the new shaded region - except that each of the two equilateral triangles will be counted twice. The area of an equilateral triangle given its side, $s,$ is $\frac{\sqrt3}4s^2.$

 Therefore, the area of the new shaded region is $4\cdot \left( \frac 16 \cdot \pi\cdot 2^2 \right) - 2 \cdot \left( 4 \cdot \frac{\sqrt3}4 \right) = \frac 83 \pi - 2\sqrt 3$. Lastly, we must subtract the area of the circle that we added earlier, $\pi$, and we get

 $\left( \frac 83 \pi - 2\sqrt 3 \right) - \pi = \boxed{\mathrm{(B)\ } \frac 53 \pi - 2\sqrt 3 }$.

 



\textbf{Solution 2 (Cheap - Out Of Time)}

Notice that answer choices (C), (D), and (E) are exactly $\pi$ more than answer choices (A), (B). This suggests that answer choices (C), (D), and (E) are meant to trick you into forgetting to subtract the circle with radius $1$ (area $\pi$) From this, (A) or (B) are probably the answer choices (since you do need to subtract $\pi$ at the end!) Since the angle measures of the sectors are $120^{\circ},$ and we are working with the areas of some equilateral triangles, it only makes sense for the answer to have a $\sqrt 3$ and not a $\sqrt 2$ in them (There's no $45^{\circ}$ angle either). So, our answer is $\boxed{\mathrm{(B)\ } \frac 53 \pi - 2\sqrt 3 }$.

 \textbf{=Video Solution 1 by SpreadTheMathLove} 


\end{enumerate}